\subsubsection{Dijkstra}
\begin{lstlisting}[style=mystyle, language=C++]
// TC: O((N + M) log N)
// usa: camino minimo desde s con pesos no negativos; requiere pesos >= 0
vector<int> dijkstra(int n, vector<vector<int>> adj[], int s){
	set<array<int,2>> st;
	vector<int> dist(n,1e18);
	dist[s]=0;
	st.insert({0,s});
	while(!st.empty()){
		auto it=*st.begin();
		int dis=it[0] , node=it[1];
		st.erase(it);
		for(auto [adjN,adjW] : adj[node]){
			if(dis+adjW < dist[adjN]){
				if(dist[adjN] != 1e18){
					st.erase({dist[adjN],adjN});
				}
				dist[adjN]=dis+adjW;
				st.insert({dist[adjN],adjN});
			}
		}
	}
	return dist;
}
// TC: O((N + M) log M)
// usa: camino minimo desde s con pesos no negativos; inserta duplicados y descarta entradas viejas al sacar
vector<int> dijkstraPQ(int n, vector<vector<int>> adj[], int s){
	priority_queue<pair<int,int>, vector<pair<int,int>>, greater<pair<int,int>>> pq;
	vector<int> dist(n,1e18);
	dist[s]=0;
	pq.push({0,s});
	while(!pq.empty()){
		auto [dis,node]=pq.top(); pq.pop();
		if(dis>dist[node]) continue;
		for(auto [adjN,adjW] : adj[node]){
			if(dis+adjW < dist[adjN]){
				dist[adjN]=dis+adjW;
				pq.push({dist[adjN],adjN});
			}
		}
	}
	return dist;
}
\end{lstlisting}
